% =============================================================================
%  3-Conclusiones.tex — Discusión, implicaciones y cierre
% =============================================================================

\section{Discusión e implicaciones}

La integración exitosa de Apimetro demuestra que la colaboración entre datos
gubernamentales (GTFS) y cartografía colaborativa (OpenStreetMap) abre un canal
inédito para la investigación urbana en América Latina.  Al empaquetar una base de
datos PostGIS compleja detrás de una API REST que entrega GeoJSON limpios, se logra
el cometido de la democratización: cualquier urbanista, geógrafo o estudiante puede,
a través de protocolos HTTP, inyectar la red de transporte completa de la ZMVM
directamente en herramientas como QGIS, Kepler.gl o Leaflet, sin necesidad de
programar ni reparar datos.

La arquitectura de Apimetro es deliberadamente replicable.  Cualquier ciudad
latinoamericana que publique \textit{feeds} GTFS y cuente con cobertura de
OpenStreetMap puede instanciar su propio motor de datos metropolitanos siguiendo
el mismo patrón: ingesta, normalización, unificación con OSM y exposición vía API
REST.  La herramienta trasciende el nicho informático para convertirse en un bien
público digital indispensable para el análisis socioespacial, el diseño de políticas
públicas de movilidad y el fortalecimiento del ecosistema de datos geoespaciales
abiertos en la región.

Se invita a la comunidad de OpenStreetMap en Latinoamérica a colaborar en la
mejora de la cobertura de datos de transporte público, particularmente en las
periferias metropolitanas donde la información gubernamental es más escasa y donde
la cartografía colaborativa tiene mayor potencial de impacto.

\subsection{Líneas de trabajo abiertas}

El desarrollo de Apimetro abre al menos tres frentes de colaboración con la
comunidad OSM:

\begin{enumerate}[nosep, leftmargin=1.5em]
  \item \textbf{Mapatones temáticos:}  Organizar sesiones de mapeo enfocadas en
    sistemas de transporte periférico (colectivos, microbuses, rutas alimentadoras)
    cuya cobertura OSM es aún incompleta en municipios conurbados como Ecatepec,
    Nezahualcóyotl y Chimalhuacán.
  \item \textbf{Validación bidireccional:}  Publicar como \textit{datasets}
    abiertos los resultados del proceso de \textit{conflation} de Apimetro, de
    modo que la comunidad pueda corregir errores detectados tanto en OSM como en
    las fuentes GTFS oficiales.
  \item \textbf{Replicabilidad regional:}  Documentar el patrón
    GTFS\,+\,OSM\,$\to$\,API REST como un modelo reproducible para otras zonas
    metropolitanas latinoamericanas ---Bogotá, São Paulo, Buenos Aires--- que
    enfrentan fragmentación similar de datos de transporte.
\end{enumerate}

Apimetro demuestra que la infraestructura de datos abiertos y la cartografía
colaborativa no son esfuerzos paralelos sino complementarios: la calidad de uno
potencia la utilidad del otro.  La convergencia entre datos gubernamentales y
comunitarios constituye, a juicio de los autores, el camino más eficiente hacia la
soberanía tecnológica en materia de información geoespacial urbana.
